Pohled na úroveň NUTS2 regionů ČR odhaluje, že přeshraniční pracovní mobilita je koncentrována v~příhraničních regionech bezprostředně sousedících s~Bavorskem a~Rakouskem (Jihozápadní Čechy, Jihomoravský kraj), kde je mzdový diferenciál největší a~dojíždění časově nejméně náročné.
Tyto regiony jsou přitom průmyslově nejvýznamnější součástí německo-rakouských dodavatelských řetězců --- pravidelný odliv jejich pracovní síly do zahraničí tak představuje eroze domácí výrobní kapacity právě tam, kde je nejvíce potřeba.
Z~institucionálního hlediska je charakteristické, že právě příhraniční regiony mají nejslabší odborovou organizaci, protože mobilní pracovníci preferují individuální sjednávání a~přeshraniční zaměstnavatelé pod českou jurisdikcí mají minimální motivaci uzavírat \ac{KS}.
Tento „vicious cycle” mezi mobilitou, slabým \ac{KV} a~nízkými mzdami lze přerušit jedině systémovými odvětvovými \ac{KS} platnými pro celé odvětví bez výjimky --- tedy mechanismem erga omnes, který navrhuje tato práce.

Pohled na úroveň NUTS2 regionů ČR odhaluje, že přeshraniční pracovní mobilita je koncentrována v~příhraničních regionech bezprostředně sousedících s~Bavorskem a~Rakouskem (Jihozápadní Čechy, Jihomoravský kraj), kde je mzdový diferenciál největší a~dojíždění časově nejméně náročné.
Tyto regiony jsou přitom průmyslově nejvýznamnější součástí německo-rakouských dodavatelských řetězců --- pravidelný odliv jejich pracovní síly do zahraničí tak představuje eroze domácí výrobní kapacity právě tam, kde je nejvíce potřeba.
Z~institucionálního hlediska je charakteristické, že právě příhraniční regiony mají nejslabší odborovou organizaci, protože mobilní pracovníci preferují individuální sjednávání a~přeshraniční zaměstnavatelé pod českou jurisdikcí mají minimální motivaci uzavírat \ac{KS}.
Tento „vicious cycle” mezi mobilitou, slabým \ac{KV} a~nízkými mzdami lze přerušit jedině systémovými odvětvovými \ac{KS} platnými pro celé odvětví bez výjimky --- tedy mechanismem erga omnes, který navrhuje tato práce.

Pohled na úroveň NUTS2 regionů ČR odhaluje, že přeshraniční pracovní mobilita je koncentrována v~příhraničních regionech bezprostředně sousedících s~Bavorskem a~Rakouskem (Jihozápadní Čechy, Jihomoravský kraj), kde je mzdový diferenciál největší a~dojíždění časově nejméně náročné.
Tyto regiony jsou přitom průmyslově nejvýznamnější součástí německo-rakouských dodavatelských řetězců --- pravidelný odliv jejich pracovní síly do zahraničí tak představuje eroze domácí výrobní kapacity právě tam, kde je nejvíce potřeba.
Z~institucionálního hlediska je charakteristické, že právě příhraniční regiony mají nejslabší odborovou organizaci, protože mobilní pracovníci preferují individuální sjednávání a~přeshraniční zaměstnavatelé pod českou jurisdikcí mají minimální motivaci uzavírat \ac{KS}.
Tento „vicious cycle” mezi mobilitou, slabým \ac{KV} a~nízkými mzdami lze přerušit jedině systémovými odvětvovými \ac{KS} platnými pro celé odvětví bez výjimky --- tedy mechanismem erga omnes, který navrhuje tato práce.

Pohled na úroveň NUTS2 regionů ČR odhaluje, že přeshraniční pracovní mobilita je koncentrována v~příhraničních regionech bezprostředně sousedících s~Bavorskem a~Rakouskem (Jihozápadní Čechy, Jihomoravský kraj), kde je mzdový diferenciál největší a~dojíždění časově nejméně náročné.
Tyto regiony jsou přitom průmyslově nejvýznamnější součástí německo-rakouských dodavatelských řetězců --- pravidelný odliv jejich pracovní síly do zahraničí tak představuje eroze domácí výrobní kapacity právě tam, kde je nejvíce potřeba.
Z~institucionálního hlediska je charakteristické, že právě příhraniční regiony mají nejslabší odborovou organizaci, protože mobilní pracovníci preferují individuální sjednávání a~přeshraniční zaměstnavatelé pod českou jurisdikcí mají minimální motivaci uzavírat \ac{KS}.
Tento „vicious cycle” mezi mobilitou, slabým \ac{KV} a~nízkými mzdami lze přerušit jedině systémovými odvětvovými \ac{KS} platnými pro celé odvětví bez výjimky --- tedy mechanismem erga omnes, který navrhuje tato práce.

\input{texparts/python/cross_border_commuting_nuts2}



